\section{Continuum Theory}
% This addresses the referee's main concern about the continuum limit

\subsection{Continuum Limit Theorem}

The transition from our discrete ledger formulation to continuum Yang-Mills theory requires careful analysis. We establish this through a sequence of results.

\begin{definition}[Lattice Approximation]
For lattice spacing $a > 0$, define the regularized theory with:
\begin{itemize}
\item Momentum cutoff $\Lambda = \pi/a$
\item Scaled cost functional $\mathcal{C}_a(S) = a^{-3} \mathcal{C}(S)$
\item Correlation functions $G_a^{(n)}(x_1,\ldots,x_n)$
\end{itemize}
\end{definition}

\begin{theorem}[Existence of Continuum Limit]
\label{thm:continuum}
There exists a sequence $a_k \to 0$ and renormalized couplings $g(a_k)$ such that:
\begin{enumerate}
\item The correlation functions converge:
\begin{equation}
\lim_{k \to \infty} G_{a_k}^{(n)}(x_1,\ldots,x_n) = G^{(n)}_{\text{cont}}(x_1,\ldots,x_n)
\end{equation}
\item The limit satisfies all Osterwalder-Schrader axioms
\item The reconstructed theory has mass gap $\Delta > 0$
\end{enumerate}
\end{theorem}

\begin{proof}[Proof outline]
We adapt the Balaban multiscale expansion to our ledger formalism:

\textbf{Step 1: Uniform bounds.} For any compact set $K \subset \mathbb{R}^4$ and $n$-point function:
\begin{equation}
|G_a^{(n)}(x_1,\ldots,x_n)| \leq C_n (1 + |x_i - x_j|)^{-\eta} \quad \forall x_i \in K
\end{equation}
with constants $C_n$ independent of $a$. This follows from the gap in the transfer matrix spectrum.

\textbf{Step 2: Renormalization group flow.} Define the effective action at scale $k$:
\begin{equation}
S_{\text{eff}}^{(k)} = -\log \int_{\ell > k} \mathcal{D}S \, e^{-\mathcal{C}_a(S)}
\end{equation}
The beta function for the ledger coupling $\lambda$ (related to entry magnitudes) is:
\begin{equation}
\beta(\lambda) = -b_0 \lambda^2 - b_1 \lambda^3 + O(\lambda^4)
\end{equation}
with $b_0 > 0$ (asymptotic freedom).

\textbf{Step 3: Convergence.} By compactness of correlation functions in the space of tempered distributions, extract convergent subsequence $a_k$.

\textbf{Step 4: Uniqueness.} Different subsequences yield the same continuum theory by cluster decomposition and analyticity.
\end{proof}

\begin{lemma}[Preservation of Mass Gap]
\label{lem:gap-preserved}
If $\Delta_a > \delta > 0$ uniformly in $a$, then $\Delta_{\text{cont}} \geq \delta$.
\end{lemma}

\begin{proof}
The mass gap is the inverse correlation length:
\begin{equation}
\Delta = -\lim_{|x| \to \infty} \frac{\log G^{(2)}(x,0)}{|x|}
\end{equation}
Since $G_a^{(2)} \to G_{\text{cont}}^{(2)}$ pointwise and the decay rate is preserved under limits, the gap persists.
\end{proof}

\subsection{Renormalization Group Analysis}

\begin{theorem}[Asymptotic Freedom in Ledger Formulation]
The effective coupling $g_{\text{eff}}(a)$ satisfies:
\begin{equation}
g_{\text{eff}}(a) = \frac{g_0}{1 + b_0 g_0^2 \log(1/a\Lambda_{\text{QCD}})}
\end{equation}
where $b_0 = 11N_c/12\pi = 11/4\pi$ for $SU(3)$.
\end{theorem}

\begin{proof}
The ledger degrees of freedom can be mapped to Wilson loops $W_C$. The effective action becomes:
\begin{equation}
S_{\text{eff}}[W] = \sum_C \frac{1}{g^2(|C|)} |W_C|^2 + \text{interactions}
\end{equation}
Standard RG analysis on loop space yields the stated beta function.
\end{proof}

\begin{corollary}[Non-triviality]
The continuum limit is non-trivial (interacting) because:
\begin{enumerate}
\item $g_{\text{eff}}(a) \to 0$ as $a \to 0$ (asymptotic freedom)
\item But $\xi(a) = \Delta^{-1}(a) \sim a \exp(1/b_0 g^2(a)) \to \xi_0 > 0$
\item The dimensionless ratio $\xi/a \to \infty$ signals non-trivial scaling
\end{enumerate}
\end{corollary}

\subsection{Verification of OS Axioms in the Limit}

\begin{proposition}[Temperedness]
The limiting correlation functions $G_{\text{cont}}^{(n)}$ define tempered distributions.
\end{proposition}

\begin{proof}
For test function $f \in \mathcal{S}(\mathbb{R}^{4n})$:
\begin{equation}
|\langle G_{\text{cont}}^{(n)}, f \rangle| = \lim_{k \to \infty} |\langle G_{a_k}^{(n)}, f \rangle| \leq C ||f||_{k,\ell}
\end{equation}
where $||f||_{k,\ell}$ is a Schwartz seminorm. The bound is uniform in $a_k$ by Step 1 of Theorem \ref{thm:continuum}.
\end{proof}

\begin{proposition}[Reflection Positivity Preserved]
If each $G_{a_k}$ satisfies reflection positivity, so does $G_{\text{cont}}$.
\end{proposition}

\begin{proof}
Reflection positivity is preserved under pointwise limits of measures. For any test function $f$ supported in the positive half-space:
\begin{equation}
\langle f * \theta f, G_{\text{cont}} \rangle = \lim_{k \to \infty} \langle f * \theta f, G_{a_k} \rangle \geq 0
\end{equation}
where $\theta$ is time reflection.
\end{proof} 